\documentclass[10pt,a4paper]{amsart}
\usepackage[margin=19mm]{geometry}
\usepackage{amsmath,amssymb,mathtools}
\usepackage[scr=rsfs,cal=euler]{mathalfa}
\usepackage{xfrac}
\usepackage[unicode,hidelinks,bookmarks=false]{hyperref}
\newcommand{\ie}{i.e.\ }
\newcommand{\cf}{cf.\ }
\newcommand{\resp}{respectively\ }
\newcommand{\Subsection}{\subsection}
\providecommand{\nobreakdash}{\nobreak\hbox{-}}
\setlength{\emergencystretch}{2em}
\multlinegap0pt
\usepackage{amssymb}
\usepackage[shortlabels]{enumitem}
\usepackage{bm}
\newtheorem{theorem}{Theorem}
\title{Covariant representations \\ of algebraic group actions \\ and applications}
\author{Yvann Gaudillot-Estrada}
\date{Source: arXiv:2603.06259v1 (CC BY 4.0)}
\begin{document}
\maketitle

\noindent\emph{Source and licence.} Covariant representations \\ of algebraic group actions \\ and applications. Version arXiv:2603.06259v1 (CC BY 4.0), distributed by arXiv under the Creative Commons Attribution 4.0 International licence, http://creativecommons.org/licenses/by/4.0/, as recorded in the arXiv licence metadata for this exact version and shown on the abstract page https://arxiv.org/abs/2603.06259v1. Full licence notice: \texttt{licenses/arxiv-2603.06259-CC-BY-4.0.txt}.\par\smallskip

\noindent\textit{Editorial scope.} Counted unit: the article's theorem classification (source0.tex lines 144--146). The article's complete original proof of it is delivered with it (source0.tex lines 227--238, one \texttt{proof} environment of 12 lines, which the article places away from the statement). No mathematics is written, repaired or re-derived: the statement and the proof are the article's own lines, in the article's own notation, and no retained line is substituted or re-flowed.  Retained uncounted as the source-local context the proof needs: lines 141--141; lines 158--160.  Nothing was omitted from any retained span, so the package carries no omission record.  No retained line contains a citation, so the wrapper carries no bibliography and no bibliography entry.  No retained line contains a figure environment, an image-inclusion command or drawing code, so no image is delivered and the package declares no asset dependency.  Editorial formatting only, in addition: an AMS \texttt{amsart} preamble that restates the article's own declarations and macros used by the retained lines, namely the environments \texttt{theorem} on separate counters, together with the standard packages those lines need.  The retained environments print in this document as Theorem 1 in order of appearance, whereas the article prints the same items with the numbering of its own full document.  The complete native source of the article as distributed in the arXiv e-print for this exact version is bundled unchanged as \texttt{sources/195851/source0.tex}.
    \begin{equation}\label{twist} T(h\cdot v) = (ghg^{-1})\cdot T(v) \qquad (v \in V, h \in G^{x}).\end{equation}

\begin{theorem}\label{equivalence}
    Assume that $G$ acts transitively on $X$. Let $x$ be a point of $X$ and let us denote by $\mathrm{ev}_x$ the functor $M \mapsto M|_x$, from $\mathrm{Rep}(G,X)$ to $\mathrm{Rep}(G^x)$. Then $\mathrm{ev}_x$ and $I_x : \mathrm{Rep}(G^x) \to \mathrm{Rep}(G,X)$ are inverse functors.
\end{theorem}

\begin{theorem}\label{classification}
    For any induction datum $(x,V)$, the covariant representation $I_x(V)$ is irreducible. Conversely, any irreducible covariant representation is isomorphic to one of that form. Moreover, for a given pair of induction data $(x,V)$ and $(x',V')$, the covariant representations $I_x(V)$ and $I_{x'}(V')$ are isomorphic if and only if $(x,V)$ is equivalent to $(x',V')$.
\end{theorem}

    \begin{proof}[Proof of Theorem \ref{classification} in the transitive case]
    
    The first two statements of the classification follow from Theorem \ref{equivalence}. It remains to check the last one.
    Let $(x,V)$ and $(x',V')$ be two induction data. Let us first assume that they are equivalent, so there exists $g \in G$ such that $x' = gx$ and a linear isomorphism $T : V\to V'$ satisfying (\ref{twist}). If $\rho$ denotes the action of $G$ on $k[G]$ by right translation, then the map
    $$\begin{array}{cclll}
        k[G] \otimes_k V & \longrightarrow & k[G] \otimes_kV' \\
        f \otimes v & \longmapsto & \rho_gf \otimes T(v) 
    \end{array}$$
    induces a $(G, X)$-isomorphism $I_x(V) \to I_{x'}(V')$.
    
    Conversely, let us assume that $I_{x'}(V')$ and $I_{x}(V)$ are isomorphic. From the foregoing, we can assume that $x = x'$ without loss of generality. Then $V$ is isomorphic to $V'$ by Theorem \ref{equivalence}.
\end{proof}

\end{document}
